\begin{thebibliography}{99}

\bibitem[ABLP26]{ABLP26}
H. Awada, M. Bolognesi, R. Laterveer and C. Pedrini,
\emph{K3 surfaces and cubic fourfolds with Abelian motive},
Ann. Mat. Pura Appl. (4), published online 2026; arXiv:2012.05969.
\href{https://doi.org/10.1007/s10231-026-01689-x}{\texttt{doi:10.1007/s10231-026-01689-x}}.

\bibitem[ACa79]{ACa79}
N. A'Campo,
\emph{Tresses, monodromie et le groupe symplectique},
Comment. Math. Helv. \textbf{54} (1979), 318--327.
\href{https://doi.org/10.1007/BF02566275}{\texttt{doi:10.1007/BF02566275}}.

\bibitem[ACGH85]{ACGH85}
E. Arbarello, M. Cornalba, P. A. Griffiths and J. Harris,
\emph{Geometry of Algebraic Curves, Volume I},
Grundlehren der Mathematischen Wissenschaften \textbf{267},
Springer, New York, 1985.

\bibitem[Alb34]{Alb34}
A. A. Albert,
\emph{On the construction of Riemann matrices. I},
Ann. of Math. (2) \textbf{35} (1934), no.~1, 1--28.
\href{https://doi.org/10.2307/1968116}{\texttt{doi:10.2307/1968116}}.

\bibitem[Alb35a]{Alb35a}
A. A. Albert,
\emph{On the construction of Riemann matrices. II},
Ann. of Math. (2) \textbf{36} (1935), no.~2, 376--394.
\href{https://doi.org/10.2307/1968578}{\texttt{doi:10.2307/1968578}}.

\bibitem[Alb35b]{Alb35b}
A. A. Albert,
\emph{Involutorial simple algebras and real Riemann matrices},
Ann. of Math. (2) \textbf{36} (1935), no.~4, 886--964.
\href{https://doi.org/10.2307/1968595}{\texttt{doi:10.2307/1968595}}.

\bibitem[Ale97]{Ale97}
H. Alexander,
\emph{Holomorphic chains and the support hypothesis conjecture},
J. Amer. Math. Soc. \textbf{10} (1997), no.~1, 123--138.
\href{https://doi.org/10.1090/S0894-0347-97-00218-X}{\texttt{doi:10.1090/S0894-0347-97-00218-X}}.

\bibitem[And92]{And92}
Y. Andr\'e,
\emph{Mumford--Tate groups of mixed Hodge structures and the theorem of the
fixed part},
Compositio Math. \textbf{82} (1992), no.~1, 1--24.
\href{http://www.numdam.org/item/CM_1992__82_1_1_0/}{\texttt{numdam.org/item/CM\_1992\_\_82\_1\_1\_0}}.

\bibitem[And96a]{And96}
Y. Andr\'e,
\emph{On the Shafarevich and Tate conjectures for hyperk\"ahler varieties},
Math. Ann. \textbf{305} (1996), no.~1, 205--248.
\href{https://doi.org/10.1007/BF01444219}{\texttt{doi:10.1007/BF01444219}}.

\bibitem[And96b]{And96b}
Y. Andr\'e,
\emph{Pour une th\'eorie inconditionnelle des motifs},
Inst. Hautes \'Etudes Sci. Publ. Math. \textbf{83} (1996), 5--49.
\href{https://doi.org/10.1007/BF02698643}{\texttt{doi:10.1007/BF02698643}}.

\bibitem[AP07]{AP07}
J. D. Achter and R. Pries,
\emph{The integral monodromy of hyperelliptic and trielliptic curves},
Math. Ann. \textbf{338} (2007), no.~1, 187--206; arXiv:math/0608038.
\href{https://doi.org/10.1007/s00208-006-0072-0}{\texttt{doi:10.1007/s00208-006-0072-0}}.

\bibitem[BB66]{BB66}
W. L. Baily, Jr. and A. Borel,
\emph{Compactification of arithmetic quotients of bounded symmetric domains},
Ann. of Math. (2) \textbf{84} (1966), no.~3, 442--528.
\href{https://doi.org/10.2307/1970457}{\texttt{doi:10.2307/1970457}}.

\bibitem[BD85]{BD85}
A. Beauville and R. Donagi,
\emph{La vari\'et\'e des droites d'une hypersurface cubique de dimension
$4$},
C. R. Acad. Sci. Paris S\'er. I Math. \textbf{301} (1985), no.~14, 703--706.

\bibitem[Bea83]{Bea83}
A. Beauville,
\emph{Vari\'et\'es K\"ahleriennes dont la premi\`ere classe de Chern est
nulle},
J. Differential Geom. \textbf{18} (1983), no.~4, 755--782.
\href{https://doi.org/10.4310/jdg/1214438181}{\texttt{doi:10.4310/jdg/1214438181}}.

\bibitem[BF03]{BF03}
R.-O. Buchweitz and H. Flenner,
\emph{A semiregularity map for modules and applications to deformations},
Compositio Math. \textbf{137} (2003), no.~2, 135--210.
\href{https://doi.org/10.1023/A:1023999012081}{\texttt{doi:10.1023/A:1023999012081}}.

\bibitem[BF08]{BF08}
R.-O. Buchweitz and H. Flenner,
\emph{The global decomposition theorem for Hochschild (co-)homology of
singular spaces via the Atiyah--Chern character},
Adv. Math. \textbf{217} (2008), no.~1, 243--281.
\href{https://doi.org/10.1016/j.aim.2007.06.013}{\texttt{doi:10.1016/j.aim.2007.06.013}}.

\bibitem[BL24]{BL24}
M. Bolognesi and R. Laterveer,
\emph{A $9$-dimensional family of K3 surfaces with finite-dimensional motive},
Rend. Circ. Mat. Palermo (2) \textbf{73} (2024), 2313--2331.
\href{https://doi.org/10.1007/s12215-024-01036-0}{\texttt{doi:10.1007/s12215-024-01036-0}}.

\bibitem[BLM23]{BLM23}
R. Bandiera, E. Lepri and M. Manetti,
\emph{$L_{\infty}$ liftings of semiregularity maps via Chern--Simons classes},
Adv. Math. \textbf{435} (2023), 109358; arXiv:2111.12985.
\href{https://doi.org/10.1016/j.aim.2023.109358}{\texttt{doi:10.1016/j.aim.2023.109358}}.

\bibitem[BMB26]{BMB26}
D. Bhattacharjee, P. Mandal and U. Bhattacharya,
\emph{Verification code and figure sources for the paper ``Algebraic Loci of
Weil Classes on Abelian Varieties''}, software archive with a Lean~4
certificate, release 5.0.1, Zenodo, 2026.
\href{https://doi.org/\zenodoCodeID}{\texttt{doi:\zenodoCodeID}}.

\bibitem[BHPV04]{BHPV04}
W. P. Barth, K. Hulek, C. A. M. Peters and A. Van de Ven,
\emph{Compact Complex Surfaces}, 2nd ed.,
Ergebnisse der Mathematik und ihrer Grenzgebiete, 3. Folge, \textbf{4},
Springer, Berlin, 2004.
\href{https://doi.org/10.1007/978-3-642-57739-0}{\texttt{doi:10.1007/978-3-642-57739-0}}.

\bibitem[BL04]{BL04}
C. Birkenhake and H. Lange,
\emph{Complex Abelian Varieties}, 2nd ed.,
Grundlehren der mathematischen Wissenschaften \textbf{302},
Springer, Berlin, 2004.
\href{https://doi.org/10.1007/978-3-662-06307-1}{\texttt{doi:10.1007/978-3-662-06307-1}}.

\bibitem[BEK14]{BEK14}
S. Bloch, H. Esnault and M. Kerz,
\emph{$p$-adic deformation of algebraic cycle classes},
Invent. Math. \textbf{195} (2014), 673--722.
\href{https://doi.org/10.1007/s00222-013-0461-4}{\texttt{doi:10.1007/s00222-013-0461-4}}.

\bibitem[Bor60]{Bor60}
A. Borel,
\emph{Density properties for certain subgroups of semi-simple groups without
compact components},
Ann. of Math. (2) \textbf{72} (1960), 179--188.
\href{https://doi.org/10.2307/1970150}{\texttt{doi:10.2307/1970150}}.

\bibitem[Bry87]{Bry87}
R. L. Bryant,
\emph{Metrics with exceptional holonomy},
Ann. of Math. (2) \textbf{126} (1987), no.~3, 525--576.

\bibitem[BS83]{BS83}
S. Bloch and V. Srinivas,
\emph{Remarks on correspondences and algebraic cycles},
Amer. J. Math. \textbf{105} (1983), no.~5, 1235--1253.
\href{https://doi.org/10.2307/2374341}{\texttt{doi:10.2307/2374341}}.

\bibitem[BS94]{BS94}
S. Bando and Y.-T. Siu,
\emph{Stable sheaves and Einstein--Hermitian metrics},
in: Geometry and Analysis on Complex Manifolds,
World Sci. Publ., River Edge, NJ, 1994, pp.~39--50.
\href{https://doi.org/10.1142/9789814350112_0002}{\texttt{doi:10.1142/9789814350112\_0002}}.

\bibitem[Blo72]{Blo72}
S. Bloch,
\emph{Semi-regularity and de Rham cohomology},
Invent. Math. \textbf{17} (1972), 51--66.
\href{https://doi.org/10.1007/BF01390023}{\texttt{doi:10.1007/BF01390023}}.

\bibitem[B\"ul20]{Bue20}
T.-H. B\"ulles,
\emph{Motives of moduli spaces on K3 surfaces and of special cubic
fourfolds},
Manuscripta Math. \textbf{161} (2020), 109--124; arXiv:1806.08284.
\href{https://doi.org/10.1007/s00229-018-1086-0}{\texttt{doi:10.1007/s00229-018-1086-0}}.

\bibitem[Bus19]{Bus19}
N. Buskin,
\emph{Every rational Hodge isometry between two K3 surfaces is algebraic},
J. Reine Angew. Math. \textbf{755} (2019), 127--150.
\href{https://doi.org/10.1515/crelle-2017-0027}{\texttt{doi:10.1515/crelle-2017-0027}}.

\bibitem[Cal05]{Cal05}
A. C\u{a}ld\u{a}raru,
\emph{The Mukai pairing, II: the Hochschild--Kostant--Rosenberg isomorphism},
Adv. Math. \textbf{194} (2005), no.~1, 34--66.
\href{https://doi.org/10.1016/j.aim.2004.05.012}{\texttt{doi:10.1016/j.aim.2004.05.012}}.

\bibitem[Cam92]{Cam92}
F. Campana,
\emph{Connexit\'e rationnelle des vari\'et\'es de Fano},
Ann. Sci. \'Ecole Norm. Sup. (4) \textbf{25} (1992), no.~5, 539--545.
\href{https://doi.org/10.24033/asens.1658}{\texttt{doi:10.24033/asens.1658}}.

\bibitem[CDK95]{CDK95}
E. Cattani, P. Deligne and A. Kaplan,
\emph{On the locus of Hodge classes},
J. Amer. Math. Soc. \textbf{8} (1995), no.~2, 483--506.
\href{https://doi.org/10.1090/S0894-0347-1995-1273413-2}{\texttt{doi:10.1090/S0894-0347-1995-1273413-2}}.

\bibitem[CKS86]{CKS86}
E. Cattani, A. Kaplan and W. Schmid,
\emph{Degeneration of Hodge structures},
Ann. of Math. (2) \textbf{123} (1986), no.~3, 457--535.
\href{https://doi.org/10.2307/1971333}{\texttt{doi:10.2307/1971333}}.

\bibitem[CLS11]{CLS11}
D. A. Cox, J. B. Little and H. K. Schenck,
\emph{Toric Varieties},
Grad. Stud. Math. \textbf{124}, Amer. Math. Soc., Providence, RI, 2011.
\href{https://doi.org/10.1090/gsm/124}{\texttt{doi:10.1090/gsm/124}}.

\bibitem[CM13]{CM13}
F. Charles and E. Markman,
\emph{The standard conjectures for holomorphic symplectic varieties
deformation equivalent to Hilbert schemes of K3 surfaces},
Compos. Math. \textbf{149} (2013), no.~3, 481--494.
\href{https://doi.org/10.1112/S0010437X12000607}{\texttt{doi:10.1112/S0010437X12000607}}.

\bibitem[CM78]{CM78}
A. Conte and J. P. Murre,
\emph{The Hodge conjecture for fourfolds admitting a covering by rational
curves},
Math. Ann. \textbf{238} (1978), 79--88.
\href{https://doi.org/10.1007/BF01351457}{\texttt{doi:10.1007/BF01351457}}.

\bibitem[CRV12]{CRV12}
D. Calaque, C. A. Rossi and M. Van den Bergh,
\emph{C\u{a}ld\u{a}raru's conjecture and Tsygan's formality},
Ann. of Math. (2) \textbf{176} (2012), no.~2, 865--923.
\href{https://doi.org/10.4007/annals.2012.176.2.4}{\texttt{doi:10.4007/annals.2012.176.2.4}}.

\bibitem[dCM02]{dCM02}
M. A. de Cataldo and L. Migliorini,
\emph{The Chow groups and the motive of the Hilbert scheme of points on a
surface},
J. Algebra \textbf{251} (2002), no.~2, 824--848.
\href{https://doi.org/10.1006/jabr.2001.9105}{\texttt{doi:10.1006/jabr.2001.9105}}.

\bibitem[Del71]{Del71}
P. Deligne,
\emph{Th\'eorie de Hodge. II},
Inst. Hautes \'Etudes Sci. Publ. Math. \textbf{40} (1971), 5--57.
\href{https://doi.org/10.1007/BF02684692}{\texttt{doi:10.1007/BF02684692}}.

\bibitem[Del72]{Del72}
P. Deligne,
\emph{La conjecture de Weil pour les surfaces $K3$},
Invent. Math. \textbf{15} (1972), 206--226.
\href{https://doi.org/10.1007/BF01404126}{\texttt{doi:10.1007/BF01404126}}.

\bibitem[Del74]{Del74}
P. Deligne,
\emph{La conjecture de Weil. I},
Inst. Hautes \'Etudes Sci. Publ. Math. \textbf{43} (1974), 273--307.
\href{https://doi.org/10.1007/BF02684373}{\texttt{doi:10.1007/BF02684373}}.

\bibitem[Del82]{Del82}
P. Deligne,
\emph{Hodge cycles on abelian varieties} (notes by J. S. Milne),
in: Hodge Cycles, Motives, and Shimura Varieties,
Lecture Notes in Math. \textbf{900}, Springer, Berlin, 1982, pp.~9--100.
\href{https://doi.org/10.1007/978-3-540-38955-2_3}%
{\texttt{doi:10.1007/978-3-540-38955-2\_3}}.

\bibitem[Dem82]{Dem82}
J.-P. Demailly,
\emph{Courants positifs extr\^emaux et conjecture de Hodge},
Invent. Math. \textbf{69} (1982), no.~3, 347--374.
\href{https://doi.org/10.1007/BF01389359}{\texttt{doi:10.1007/BF01389359}}.

\bibitem[Dem12]{Dem12}
J.-P. Demailly,
\emph{Complex Analytic and Differential Geometry},
OpenContent book, Institut Fourier, Grenoble, 2012,
\href{https://www-fourier.univ-grenoble-alpes.fr/~demailly/manuscripts/agbook.pdf}{\texttt{www-fourier.univ-grenoble-alpes.fr/\~{}demailly/manuscripts/agbook.pdf}}.

\bibitem[DM82]{DM82}
P. Deligne and J. S. Milne,
\emph{Tannakian categories},
in: Hodge Cycles, Motives, and Shimura Varieties,
Lecture Notes in Math. \textbf{900}, Springer, Berlin, 1982, pp.~101--228.
\href{https://doi.org/10.1007/978-3-540-38955-2_4}%
{\texttt{doi:10.1007/978-3-540-38955-2\_4}}.

\bibitem[DM86]{DM86}
P. Deligne and G. D. Mostow,
\emph{Monodromy of hypergeometric functions and non-lattice integral
monodromy},
Inst. Hautes \'Etudes Sci. Publ. Math. \textbf{63} (1986), 5--89.

\bibitem[EGL01]{EGL01}
G. Ellingsrud, L. G\"ottsche and M. Lehn,
\emph{On the cobordism class of the Hilbert scheme of a surface},
J. Algebraic Geom. \textbf{10} (2001), no.~1, 81--100.

\bibitem[Fed69]{Fed69}
H. Federer,
\emph{Geometric Measure Theory},
Grundlehren der mathematischen Wissenschaften \textbf{153},
Springer, Berlin, 1969.

\bibitem[Fed74]{Fed74}
H. Federer,
\emph{Real flat chains, cochains and variational problems},
Indiana Univ. Math. J. \textbf{24} (1974/75), no.~4, 351--407.
\href{https://doi.org/10.1512/iumj.1975.24.24031}{\texttt{doi:10.1512/iumj.1975.24.24031}}.

\bibitem[FF60]{FF60}
H. Federer and W. H. Fleming,
\emph{Normal and integral currents},
Ann. of Math. (2) \textbf{72} (1960), no.~3, 458--520.
\href{https://doi.org/10.2307/1970227}{\texttt{doi:10.2307/1970227}}.

\bibitem[FFZ21]{FFZ21}
S. Floccari, L. Fu and Z. Zhang,
\emph{On the motive of O'Grady's ten-dimensional hyper-K\"ahler varieties},
Commun. Contemp. Math. \textbf{23} (2021), no.~4, paper no.~2050034,
50~pp.
\href{https://doi.org/10.1142/S0219199720500340}{\texttt{doi:10.1142/S0219199720500340}}.

\bibitem[FK26]{FK26}
D. Favero and T. L. Kelly,
\emph{The Chern character of a coherent sheaf on a smooth projective
hypersurface},
preprint, 2026.
\href{https://arxiv.org/abs/2609.12759}{\texttt{arXiv:2609.12759}}.

\bibitem[Flo23]{Flo23}
S. Floccari,
\emph{On the motive of O'Grady's six-dimensional hyper-K\"ahler varieties},
\'Epijournal G\'eom. Alg\'ebrique \textbf{7} (2023), Article no.~4;
arXiv:2203.16257.
\href{https://doi.org/10.46298/epiga.2022.9758}{\texttt{doi:10.46298/epiga.2022.9758}}.

\bibitem[Flo25]{Flo25}
S. Floccari,
\emph{K3 surfaces associated with varieties of generalized Kummer type},
Geom. Topol. \textbf{30} (2026), no.~3, 1129--1154; arXiv:2501.02315.
\href{https://doi.org/10.2140/gt.2026.30.1129}{\texttt{doi:10.2140/gt.2026.30.1129}}.

\bibitem[FM98]{FM98}
B. Fantechi and M. Manetti,
\emph{Obstruction calculus for functors of Artin rings, I},
J. Algebra \textbf{202} (1998), no.~2, 541--576.
\href{https://doi.org/10.1006/jabr.1997.7239}{\texttt{doi:10.1006/jabr.1997.7239}}.

\bibitem[FF26]{FF26}
S. Floccari and L. Fu,
\emph{The Hodge conjecture for Weil fourfolds with discriminant $1$ via
singular OG6-varieties},
J. Math. Pures Appl. (9) \textbf{210} (2026), paper no.~103876.
\href{https://doi.org/10.1016/j.matpur.2026.103876}{\texttt{doi:10.1016/j.matpur.2026.103876}}.

\bibitem[Flo24]{Flo24}
S. Floccari,
\emph{Sixfolds of generalized Kummer type and K3 surfaces},
Compos. Math. \textbf{160} (2024), no.~2, 388--410.
\href{https://doi.org/10.1112/S0010437X23007625}{\texttt{doi:10.1112/S0010437X23007625}}.

\bibitem[Fox53]{Fox53}
R. H. Fox,
\emph{Free differential calculus. I. Derivation in the free group ring},
Ann. of Math. (2) \textbf{57} (1953), no.~3, 547--560.

\bibitem[FTV19]{FTV19}
L. Fu, Z. Tian and C. Vial,
\emph{Motivic hyper-K\"ahler resolution conjecture, I: generalized Kummer
varieties},
Geom. Topol. \textbf{23} (2019), no.~1, 427--492; arXiv:1608.04968.
\href{https://doi.org/10.2140/gt.2019.23.427}{\texttt{doi:10.2140/gt.2019.23.427}}.

\bibitem[Fuj87]{Fuj87}
A. Fujiki,
\emph{On the de Rham cohomology group of a compact K\"ahler symplectic
manifold},
in: Algebraic Geometry, Sendai, 1985, Adv. Stud. Pure Math. \textbf{10}
(1987), 105--165.
\href{https://doi.org/10.2969/aspm/01010105}{\texttt{doi:10.2969/aspm/01010105}}.

\bibitem[Ful98]{Ful98}
W. Fulton,
\emph{Intersection Theory}, 2nd ed.,
Ergebnisse der Mathematik und ihrer Grenzgebiete (3) \textbf{2},
Springer, New York, 1998.
\href{https://doi.org/10.1007/978-1-4612-1700-8}{\texttt{doi:10.1007/978-1-4612-1700-8}}.

\bibitem[FH91]{FH91}
W. Fulton and J. Harris,
\emph{Representation Theory: A First Course},
Graduate Texts in Mathematics \textbf{129}, Springer, New York, 1991.

\bibitem[FV25]{FV25}
S. Floccari and M. Varesco,
\emph{Algebraic cycles on hyper-K\"ahler varieties of generalized Kummer
type},
Math. Ann. \textbf{391} (2025), 4443--4453; arXiv:2308.04865.
\href{https://doi.org/10.1007/s00208-024-03027-z}{\texttt{doi:10.1007/s00208-024-03027-z}}.

\bibitem[Gal00]{Gal00}
F. Galluzzi,
\emph{Abelian fourfold of Mumford-type and Kuga--Satake varieties},
Indag. Math. (N.S.) \textbf{11} (2000), no.~4, 547--560.
\href{https://doi.org/10.1016/S0019-3577(00)80024-5}{\texttt{doi:10.1016/S0019-3577(00)80024-5}}.

\bibitem[GDHL09]{GDHL09}
G. Gonz\'alez-Diez, R. A. Hidalgo and M. Leyton,
\emph{Generalized Fermat curves},
J. Algebra \textbf{321} (2009), no.~6, 1643--1660.

\bibitem[GH78]{GH78}
P. Griffiths and J. Harris,
\emph{Principles of Algebraic Geometry},
Wiley, New York, 1978.

\bibitem[GM88]{GM88}
W. M. Goldman and J. J. Millson,
\emph{The deformation theory of representations of fundamental groups of
compact K\"ahler manifolds},
Inst. Hautes \'Etudes Sci. Publ. Math. \textbf{67} (1988), 43--96.
\href{https://doi.org/10.1007/BF02699127}{\texttt{doi:10.1007/BF02699127}}.

\bibitem[GP98]{GP98}
M. Gross and S. Popescu,
\emph{Equations of $(1,d)$-polarized abelian surfaces},
Math. Ann. \textbf{310} (1998), no.~2, 333--377.
\href{https://doi.org/10.1007/s002080050151}{\texttt{doi:10.1007/s002080050151}}.

\bibitem[Gri10]{Gri10}
J. Grivaux,
\emph{Chern classes in Deligne cohomology for coherent analytic sheaves},
Math. Ann. \textbf{347} (2010), no.~2, 249--284.
\href{https://doi.org/10.1007/s00208-009-0430-9}{\texttt{doi:10.1007/s00208-009-0430-9}}.

\bibitem[Gro61]{Gro61}
A. Grothendieck,
\emph{Techniques de construction et th\'eor\`emes d'existence en g\'eom\'etrie
alg\'ebrique IV: les sch\'emas de Hilbert},
S\'eminaire Bourbaki, Vol.~6 (1960/61), Exp. No.~221, 249--276;
reprinted by Soc. Math. France, Paris, 1995.

\bibitem[GW09]{GW09}
R. Goodman and N. R. Wallach,
\emph{Symmetry, Representations, and Invariants},
Graduate Texts in Mathematics \textbf{255}, Springer, New York, 2009.
\href{https://doi.org/10.1007/978-0-387-79852-3}{\texttt{doi:10.1007/978-0-387-79852-3}}.

\bibitem[Har77]{Har77}
R. Hartshorne,
\emph{Algebraic geometry},
Grad. Texts in Math. \textbf{52}, Springer, New York, 1977.
\href{https://doi.org/10.1007/978-1-4757-3849-0}{\texttt{doi:10.1007/978-1-4757-3849-0}}.

\bibitem[Hir64]{Hir64}
H. Hironaka,
\emph{Resolution of singularities of an algebraic variety over a field of
characteristic zero. I, II},
Ann. of Math. (2) \textbf{79} (1964), no.~1, 109--203; no.~2, 205--326.
\href{https://doi.org/10.2307/1970486}{\texttt{doi:10.2307/1970486}},
\href{https://doi.org/10.2307/1970547}{\texttt{doi:10.2307/1970547}}.

\bibitem[HK74]{HK74}
F. R. Harvey and A. W. Knapp,
\emph{Positive $(p,p)$-forms, Wirtinger's inequality and currents},
in: Value-Distribution Theory, Part A (Tulane Univ., 1972--73),
Dekker, New York, 1974, pp.~43--62.

\bibitem[HL82]{HL82}
R. Harvey and H. B. Lawson, Jr.,
\emph{Calibrated geometries},
Acta Math. \textbf{148} (1982), 47--157.
\href{https://doi.org/10.1007/BF02392726}{\texttt{doi:10.1007/BF02392726}}.

\bibitem[HS74]{HS74}
R. Harvey and B. Shiffman,
\emph{A characterization of holomorphic chains},
Ann. of Math. (2) \textbf{99} (1974), no.~3, 553--587.
\href{https://doi.org/10.2307/1971062}{\texttt{doi:10.2307/1971062}}.

\bibitem[HL09]{HL09}
F. R. Harvey and H. B. Lawson, Jr.,
\emph{Duality of positive currents and plurisubharmonic functions in
calibrated geometry},
Amer. J. Math. \textbf{131} (2009), no.~5, 1211--1240.
\href{https://arxiv.org/abs/0710.3921}{\texttt{arXiv:0710.3921}}.

\bibitem[HM73]{HM73}
G. Horrocks and D. Mumford,
\emph{A rank $2$ vector bundle on $\mathbf{P}^{4}$ with $15{,}000$ symmetries},
Topology \textbf{12} (1973), 63--81.
\href{https://doi.org/10.1016/0040-9383(73)90022-0}{\texttt{doi:10.1016/0040-9383(73)90022-0}}.

\bibitem[Hua21]{Hua21}
S. Huang,
\emph{A note on a question of Markman},
J. Pure Appl. Algebra \textbf{225} (2021), no.~9, 106673.
\href{https://doi.org/10.1016/j.jpaa.2021.106673}{\texttt{doi:10.1016/j.jpaa.2021.106673}}.

\bibitem[Huy16]{Huy16}
D. Huybrechts,
\emph{Lectures on K3 Surfaces},
Cambridge Studies in Advanced Mathematics \textbf{158},
Cambridge University Press, Cambridge, 2016.
\href{https://doi.org/10.1017/CBO9781316594193}{\texttt{doi:10.1017/CBO9781316594193}}.

\bibitem[ILP22]{ILP22}
C. Ingalls, A. Logan and O. Patashnick,
\emph{Explicit coverings of families of elliptic surfaces by squares of
curves},
Math. Z. \textbf{302} (2022), 1191--1238; arXiv:2009.07807.
\href{https://doi.org/10.1007/s00209-022-03090-9}{\texttt{doi:10.1007/s00209-022-03090-9}}.

\bibitem[Jan92]{Jan92}
U. Jannsen,
\emph{Motives, numerical equivalence, and semi-simplicity},
Invent. Math. \textbf{107} (1992), no.~3, 447--452.
\href{https://doi.org/10.1007/BF01231898}{\texttt{doi:10.1007/BF01231898}}.

\bibitem[Jum26]{Jum26}
R. Jumagulov,
\emph{The Hodge conjecture for Fermat fourfolds of odd degree at most
$199$},
preprint, 2026.
\href{https://arxiv.org/abs/2608.18134}{\texttt{arXiv:2608.18134}}.

\bibitem[Kho77]{Kho77}
A. G. Khovanskii,
\emph{Newton polyhedra and toroidal varieties},
Funct. Anal. Appl. \textbf{11} (1977), no.~4, 289--296.
\href{https://doi.org/10.1007/BF01077143}{\texttt{doi:10.1007/BF01077143}}.

\bibitem[Kin71]{Kin71}
J. R. King,
\emph{The currents defined by analytic varieties},
Acta Math. \textbf{127} (1971), 185--220.
\href{https://doi.org/10.1007/BF02392053}{\texttt{doi:10.1007/BF02392053}}.

\bibitem[Kle68]{Kle68}
S. L. Kleiman,
\emph{Algebraic cycles and the Weil conjectures},
in: Dix expos\'es sur la cohomologie des sch\'emas,
North-Holland, Amsterdam, 1968, pp.~359--386.

\bibitem[Kle94]{Kle94}
S. L. Kleiman,
\emph{The standard conjectures},
in: Motives (Seattle, WA, 1991),
Proc. Sympos. Pure Math. \textbf{55}, Part 1,
Amer. Math. Soc., Providence, RI, 1994, pp.~3--20.
\href{https://doi.org/10.1090/pspum/055.1/1265519}{\texttt{doi:10.1090/pspum/055.1/1265519}}.

\bibitem[KMM92]{KMM92}
J. Koll\'ar, Y. Miyaoka and S. Mori,
\emph{Rationally connected varieties},
J. Algebraic Geom. \textbf{1} (1992), no.~3, 429--448.

\bibitem[Koi04]{Koi04}
K. Koike,
\emph{Algebraicity of some Weil Hodge classes},
Canad. Math. Bull. \textbf{47} (2004), no.~4, 566--572.
\href{https://arxiv.org/abs/math/0211304}{\texttt{arXiv:math/0211304}}.

\bibitem[KOU23]{KOU}
B. Klingler, A. Otwinowska and D. Urbanik,
\emph{On the fields of definition of Hodge loci},
Ann. Sci. \'Ec. Norm. Sup\'er. (4) \textbf{56} (2023), no.~4, 1299--1312.
\href{https://doi.org/10.24033/asens.2555}{\texttt{doi:10.24033/asens.2555}}.

\bibitem[KS67]{KS67}
M. Kuga and I. Satake,
\emph{Abelian varieties attached to polarized $K3$-surfaces},
Math. Ann. \textbf{169} (1967), 239--242.
\href{https://doi.org/10.1007/BF01399540}{\texttt{doi:10.1007/BF01399540}}.

\bibitem[Kur65]{Kur65}
M. Kuranishi,
\emph{New proof for the existence of locally complete families of complex
structures},
in: Proc. Conf. Complex Analysis (Minneapolis, 1964),
Springer, Berlin, 1965, pp.~142--154.
\href{https://doi.org/10.1007/978-3-642-48016-4_13}{\texttt{doi:10.1007/978-3-642-48016-4\_13}}.

\bibitem[Li26]{Li26}
N. Li,
\emph{The Tate conjecture for powers of abelian fourfolds and the Hodge
conjecture for powers of CM fourfolds},
preprint, 2026.
\href{https://arxiv.org/abs/2609.27916}{\texttt{arXiv:2609.27916}}.

\bibitem[Lie06]{Lie06}
M. Lieblich,
\emph{Moduli of complexes on a proper morphism},
J. Algebraic Geom. \textbf{15} (2006), no.~1, 175--206.
\href{https://doi.org/10.1090/S1056-3911-05-00418-2}%
{\texttt{doi:10.1090/S1056-3911-05-00418-2}}.

\bibitem[Lie68]{Lie68}
D. I. Lieberman,
\emph{Numerical and homological equivalence of algebraic cycles on Hodge
manifolds},
Amer. J. Math. \textbf{90} (1968), no.~2, 366--374.
\href{https://doi.org/10.2307/2373533}{\texttt{doi:10.2307/2373533}}.

\bibitem[Loo97]{Loo97}
E. Looijenga,
\emph{Prym representations of mapping class groups},
Geom. Dedicata \textbf{64} (1997), 69--83.
\href{https://doi.org/10.1023/A:1004909416648}{\texttt{doi:10.1023/A:1004909416648}}.

\bibitem[Mac62]{Mac62}
I. G. Macdonald,
\emph{Symmetric products of an algebraic curve},
Topology \textbf{1} (1962), 319--343.

\bibitem[Mar24]{Mar24}
E. Markman,
\emph{Rational Hodge isometries of hyper-K\"ahler varieties of K3$^{[n]}$
type are algebraic},
Compos. Math. \textbf{160} (2024), no.~6, 1261--1303.
\href{https://doi.org/10.1112/S0010437X24007048}{\texttt{doi:10.1112/S0010437X24007048}}.

\bibitem[Mar25a]{Mar25a}
E. Markman,
\emph{Cycles on abelian $2n$-folds of Weil type from secant sheaves on abelian
$n$-folds}, preprint, arXiv:2502.03415 (2025).
\href{https://doi.org/10.48550/arXiv.2502.03415}{\texttt{doi:10.48550/arXiv.2502.03415}}.

\bibitem[Mar25b]{Mar25b}
E. Markman,
\emph{Secant sheaves and Weil classes on abelian varieties},
preprint, arXiv:2509.23403 (2025).
\href{https://doi.org/10.48550/arXiv.2509.23403}{\texttt{doi:10.48550/arXiv.2509.23403}}.

\bibitem[Mar25c]{Mar25c}
E. Markman,
\emph{Secant sheaves on abelian $n$-folds with real multiplication and Weil
classes on abelian $2n$-folds with complex multiplication},
preprint, 2025.
\href{https://arxiv.org/abs/2509.23079}{\texttt{arXiv:2509.23079}}.

\bibitem[Mat86]{Mat86}
H. Matsumura,
\emph{Commutative Ring Theory},
Cambridge Studies in Advanced Mathematics \textbf{8},
Cambridge University Press, Cambridge, 1986.
\href{https://doi.org/10.1017/CBO9781139171762}{\texttt{doi:10.1017/CBO9781139171762}}.

\bibitem[Mil20]{Mil20}
J. S. Milne,
\emph{Hodge classes on abelian varieties},
preprint, 2020.
\href{https://arxiv.org/abs/2010.08857}{\texttt{arXiv:2010.08857}}.

\bibitem[Mil22]{Mil22}
J. S. Milne,
\emph{The Tate and standard conjectures for certain abelian varieties},
preprint, arXiv:2112.12815, version 2 (2022).
\href{https://doi.org/10.48550/arXiv.2112.12815}{\texttt{doi:10.48550/arXiv.2112.12815}}.

\bibitem[MMMV26]{MMMV26}
L. Martelotte, M. Miranda, H. Movasati and R. Villaflor,
\emph{Lengths of Hodge characters in Fermat varieties},
preprint, 2026.
\href{https://arxiv.org/abs/2609.27301}{\texttt{arXiv:2609.27301}}.

\bibitem[Mor84]{Mor84}
D. R. Morrison,
\emph{On K3 surfaces with large Picard number},
Invent. Math. \textbf{75} (1984), no.~1, 105--121.
\href{https://doi.org/10.1007/BF01403093}{\texttt{doi:10.1007/BF01403093}}.

\bibitem[Mor19]{Mor19}
M. Morrow,
\emph{A variational Tate conjecture in crystalline cohomology},
J. Eur. Math. Soc. \textbf{21} (2019), no.~11, 3467--3511.
\href{https://doi.org/10.4171/JEMS/907}{\texttt{doi:10.4171/JEMS/907}}.

\bibitem[Mor85]{Mor85}
D. R. Morrison,
\emph{The Kuga--Satake variety of an abelian surface},
J. Algebra \textbf{92} (1985), 454--476.
\href{https://doi.org/10.1016/0021-8693(85)90134-6}%
{\texttt{doi:10.1016/0021-8693(85)90134-6}}.

\bibitem[Muk78]{Muk78}
S. Mukai,
\emph{Semi-homogeneous vector bundles on an abelian variety},
J. Math. Kyoto Univ. \textbf{18} (1978), no.~2, 239--272.
\href{https://doi.org/10.1215/kjm/1250522574}{\texttt{doi:10.1215/kjm/1250522574}}.

\bibitem[Muk81]{Muk81}
S. Mukai,
\emph{Duality between $D(X)$ and $D(\hat{X})$ with its application to Picard
sheaves},
Nagoya Math. J. \textbf{81} (1981), 153--175.
\href{https://doi.org/10.1017/S002776300001922X}{\texttt{doi:10.1017/S002776300001922X}}.

\bibitem[Mum69]{Mum69}
D. Mumford,
\emph{A note of Shimura's paper ``Discontinuous groups and abelian
varieties''},
Math. Ann. \textbf{181} (1969), 345--351.
\href{https://doi.org/10.1007/BF01350672}{\texttt{doi:10.1007/BF01350672}}.

\bibitem[Mur90]{Mur90}
J. P. Murre,
\emph{On the motive of an algebraic surface},
J. Reine Angew. Math. \textbf{409} (1990), 190--204.
\href{https://doi.org/10.1515/crll.1990.409.190}{\texttt{doi:10.1515/crll.1990.409.190}}.

\bibitem[MZ95]{MZ95}
B. Moonen and Yu. G. Zarhin,
\emph{Hodge classes and Tate classes on simple abelian fourfolds},
Duke Math. J. \textbf{77} (1995), no.~3, 553--581.
\href{https://doi.org/10.1215/S0012-7094-95-07717-5}{\texttt{doi:10.1215/S0012-7094-95-07717-5}}.

\bibitem[MZ99]{MZ99}
B. Moonen and Yu. G. Zarhin,
\emph{Hodge classes on abelian varieties of low dimension},
Math. Ann. \textbf{315} (1999), no.~4, 711--733.
\href{https://doi.org/10.1007/s002080050333}{\texttt{doi:10.1007/s002080050333}}.

\bibitem[OAI26a]{OAI26a}
OpenAI,
\emph{Abelian covers, Gale correspondences and the Hodge conjecture for powers},
OpenAI Math Release preprint, September 30, 2026, not refereed.
\url{https://github.com/openai/math}.

\bibitem[OAI26b]{OAI26b}
OpenAI,
\emph{Algebraic Kuga--Satake correspondences and Hodge conjectures on a K3 quadratic locus},
OpenAI Math Release preprint, September 30, 2026, not refereed.
\url{https://github.com/openai/math}.

\bibitem[OAI26c]{OAI26c}
OpenAI,
\emph{Algebraicity of Kuga--Satake correspondences for K3 surfaces},
OpenAI Math Release preprint, October 3, 2026, not refereed.
\url{https://github.com/openai/math}.

\bibitem[OAI26d]{OAI26d}
OpenAI,
\emph{Algebraicity of Weil classes on split abelian eightfolds},
OpenAI Math Release preprint, September 18, 2026, not refereed.
\url{https://github.com/openai/math}.

\bibitem[OAI26e]{OAI26e}
OpenAI,
\emph{The rational Hodge conjecture for CM abelian varieties},
OpenAI Math Release preprint, September 30, 2026, not refereed.
\url{https://github.com/openai/math}.

\bibitem[OAI26f]{OAI26f}
OpenAI,
\emph{The rational Hodge conjecture for products of K3 surfaces},
OpenAI Math Release preprint, October 4, 2026, not refereed.
\url{https://github.com/openai/math}.

\bibitem[OAI26g]{OAI26g}
OpenAI,
\emph{Weil classes and Hodge classes on abelian powers},
OpenAI Math Release preprint, September 30, 2026, not refereed.
\url{https://github.com/openai/math}.

\bibitem[O'G26]{OG26}
K. G. O'Grady,
\emph{HK manifolds of type K3$^{[a^{2}+1]}$ as moduli spaces of projective
bundles on HK manifolds of type K3$^{[2]}$},
preprint, 2026.
\href{https://arxiv.org/abs/2606.03775}{\texttt{arXiv:2606.03775}}.

\bibitem[OMe00]{OMe00}
O. T. O'Meara,
\emph{Introduction to Quadratic Forms},
Classics in Mathematics, Springer, Berlin, 2000 (reprint of the 1973 edition).
\href{https://doi.org/10.1007/978-3-642-62031-7}{\texttt{doi:10.1007/978-3-642-62031-7}}.

\bibitem[Ped26]{Ped26}
C. Pedrini,
\emph{The Chow motive of LSV hyper-K\"ahler tenfolds},
preprint, 2026.
\href{https://arxiv.org/abs/2603.10894}{\texttt{arXiv:2603.10894}}.

\bibitem[Per26]{Per26}
A. Perry,
\emph{The semiregularity theorem for equivariant noncommutative varieties},
preprint, arXiv:2604.00511 (2026).
\href{https://doi.org/10.48550/arXiv.2604.00511}{\texttt{doi:10.48550/arXiv.2604.00511}}.

\bibitem[PR94]{PR94}
V. P. Platonov and A. S. Rapinchuk,
\emph{Algebraic Groups and Number Theory},
Pure and Applied Mathematics \textbf{139},
Academic Press, Boston, MA, 1994.

\bibitem[Pri24]{Pri24}
J. P. Pridham,
\emph{Semiregularity as a consequence of Goodwillie's theorem},
Forum Math. Sigma \textbf{12} (2024), e126. arXiv:1208.3111.
\href{https://doi.org/10.1017/fms.2024.132}{\texttt{doi:10.1017/fms.2024.132}}.

\bibitem[PST21]{PST}
J. Pila, A. N. Shankar and J. Tsimerman,
\emph{Canonical heights on Shimura varieties and the Andr\'e--Oort
conjecture}, with an appendix by H. Esnault and M. Groechenig,
preprint, arXiv:2109.08788 (2021).
\href{https://doi.org/10.48550/arXiv.2109.08788}{\texttt{doi:10.48550/arXiv.2109.08788}}.

\bibitem[PZ25]{PZ25}
D. Patel and Y. Zhang,
\emph{Algebraicity of Hodge classes on some generalized Prym varieties},
preprint, 2025.
\href{https://arxiv.org/abs/2506.13729}{\texttt{arXiv:2506.13729}}.

\bibitem[RM08]{RM08}
J. J. Ram\'on Mar\'i,
\emph{On the Hodge conjecture for products of certain surfaces},
Collect. Math. \textbf{59} (2008), no.~1, 1--26.
\href{https://doi.org/10.1007/BF03191179}{\texttt{doi:10.1007/BF03191179}}.

\bibitem[Sat66]{Sat66}
I. Satake,
\emph{Clifford algebras and families of abelian varieties},
Nagoya Math. J. \textbf{27} (1966), 435--446.
\href{https://doi.org/10.1017/S0027763000026295}{\texttt{doi:10.1017/S0027763000026295}}.

\bibitem[Sch10]{Sch10}
U. Schlickewei,
\emph{The Hodge conjecture for self-products of certain K3 surfaces},
J. Algebra \textbf{324} (2010), no.~3, 507--529.
\href{https://doi.org/10.1016/j.jalgebra.2010.04.010}{\texttt{doi:10.1016/j.jalgebra.2010.04.010}}.

\bibitem[Sch73]{Sch73}
W. Schmid,
\emph{Variation of Hodge structure: the singularities of the period mapping},
Invent. Math. \textbf{22} (1973), 211--319.
\href{https://doi.org/10.1007/BF01389674}{\texttt{doi:10.1007/BF01389674}}.

\bibitem[Sch85]{Sch85}
W. Scharlau,
\emph{Quadratic and Hermitian Forms},
Grundlehren der mathematischen Wissenschaften \textbf{270},
Springer, Berlin, 1985.
\href{https://doi.org/10.1007/978-3-642-69971-9}{\texttt{doi:10.1007/978-3-642-69971-9}}.

\bibitem[Sch88]{Sch88}
C. Schoen,
\emph{Hodge classes on self-products of a variety with an automorphism},
Compositio Math. \textbf{65} (1988), no.~1, 3--32.
\href{http://www.numdam.org/item?id=CM_1988__65_1_3_0}{\texttt{numdam.org/item?id=CM\_1988\_\_65\_1\_3\_0}}.

\bibitem[Sch98]{Sch98}
C. Schoen,
\emph{Addendum to: Hodge classes on self-products of a variety with an
automorphism},
Compositio Math. \textbf{114} (1998), no.~3, 329--336.
\href{https://doi.org/10.1023/A:1000566205021}{\texttt{doi:10.1023/A:1000566205021}}.

\bibitem[Ser00]{Ser00}
J.-P. Serre,
\emph{Local Algebra},
Springer Monographs in Mathematics,
Springer, Berlin, 2000.
\href{https://doi.org/10.1007/978-3-662-04203-8}{\texttt{doi:10.1007/978-3-662-04203-8}}.

\bibitem[SGA\,$4\frac{1}{2}$]{SGA4h}
P. Deligne,
\emph{Cohomologie \'etale} (SGA $4\frac{1}{2}$), with the expos\'e
[Cycle], \emph{La classe de cohomologie associ\'ee \`a un cycle}, by
A. Grothendieck, r\'edig\'e par P. Deligne, pp.~129--153,
Lecture Notes in Math. \textbf{569}, Springer, Berlin, 1977.
\href{https://doi.org/10.1007/BFb0091516}{\texttt{doi:10.1007/BFb0091516}}.

\bibitem[Shi79]{Shi79b}
T. Shioda,
\emph{The Hodge conjecture for Fermat varieties},
Math. Ann. \textbf{245} (1979), 175--184.
\href{https://doi.org/10.1007/BF01428804}{\texttt{doi:10.1007/BF01428804}}.

\bibitem[Shi86]{Shi86}
T. Shioda,
\emph{An explicit algorithm for computing the Picard number of certain
algebraic surfaces},
Amer. J. Math. \textbf{108} (1986), no.~2, 415--432.
\href{https://doi.org/10.2307/2374678}{\texttt{doi:10.2307/2374678}}.

\bibitem[SK79]{SK79}
T. Shioda and T. Katsura,
\emph{On Fermat varieties},
T\^ohoku Math. J. (2) \textbf{31} (1979), no.~1, 97--115.
\href{https://doi.org/10.2748/tmj/1178229881}{\texttt{doi:10.2748/tmj/1178229881}}.

\bibitem[Stacks]{Stacks}
The Stacks project authors,
\emph{The Stacks project},
\href{https://stacks.math.columbia.edu}{\texttt{https://stacks.math.columbia.edu}},
2026.

\bibitem[SZ11]{SZ11}
M. Sheng and K. Zuo,
\emph{Deuring's mass formula of a Mumford family},
Trans. Amer. Math. Soc. \textbf{368} (2016), no.~1, 169--207.
\href{https://doi.org/10.1090/S0002-9947-2015-06312-0}%
{\texttt{doi:10.1090/S0002-9947-2015-06312-0}}.
\href{https://arxiv.org/abs/1111.5983}{\texttt{arXiv:1111.5983}}.

\bibitem[Tan03]{Tan03}
S. G. Tankeev,
\emph{On the standard conjecture for complex abelian schemes over smooth
projective curves},
Izv. Math. \textbf{67} (2003), no.~3, 597--635.
\href{https://doi.org/10.1070/IM2003v067n03ABEH000439}{\texttt{doi:10.1070/IM2003v067n03ABEH000439}}.

\bibitem[Tat66]{Tat66}
J. Tate,
\emph{Endomorphisms of abelian varieties over finite fields},
Invent. Math. \textbf{2} (1966), 134--144.
\href{https://doi.org/10.1007/BF01404549}{\texttt{doi:10.1007/BF01404549}}.

\bibitem[Ter88]{Ter88}
T. Terasoma,
\emph{Complete intersections of hypersurfaces: the Fermat case and the
quadric case},
Japan. J. Math. (N.S.) \textbf{14} (1988), no.~2, 309--384.
\href{https://doi.org/10.4099/math1924.14.309}{\texttt{doi:10.4099/math1924.14.309}}.

\bibitem[Tod09]{Tod09}
Y. Toda,
\emph{Deformations and Fourier--Mukai transforms},
J. Differential Geom. \textbf{81} (2009), no.~1, 197--224.
\href{https://doi.org/10.4310/jdg/1228400631}{\texttt{doi:10.4310/jdg/1228400631}}.

\bibitem[Tot26]{Tot26}
B. Totaro,
\emph{On the Hodge conjecture},
guest post on the weblog of T. Tao, 11 September 2026.
\href{https://terrytao.wordpress.com/2026/09/11/on-the-hodge-conjecture/}%
{\texttt{terrytao.wordpress.com/2026/09/11/on-the-hodge-conjecture}}.

\bibitem[Tsi18]{Tsi18}
J. Tsimerman,
\emph{The Andr\'e--Oort conjecture for $\mathcal{A}_{g}$},
Ann. of Math. (2) \textbf{187} (2018), no.~2, 379--390.
\href{https://doi.org/10.4007/annals.2018.187.2.2}{\texttt{doi:10.4007/annals.2018.187.2.2}}.

\bibitem[TY21]{TY21}
J.-H. Teh and C.-J. Yang,
\emph{Real rectifiable currents, holomorphic chains and algebraic cycles},
Complex Manifolds \textbf{8} (2021), no.~1, 274--285.
\href{https://doi.org/10.1515/coma-2020-0119}{\texttt{doi:10.1515/coma-2020-0119}}.

\bibitem[Var22]{Var22}
M. Varesco,
\emph{The Hodge conjecture for powers of K3 surfaces of Picard number $16$},
Michigan Math. J., advance publication (2024), 1--32; arXiv:2203.09778.
\href{https://doi.org/10.1307/mmj/20236349}{\texttt{doi:10.1307/mmj/20236349}}.

\bibitem[Var23]{Var23}
M. Varesco,
\emph{Hodge similarities, algebraic classes, and Kuga--Satake varieties},
Math. Z. \textbf{305} (2023), paper no.~69.
\href{https://doi.org/10.1007/s00209-023-03390-8}{\texttt{doi:10.1007/s00209-023-03390-8}}.

\bibitem[vG00]{vG00}
B. van Geemen,
\emph{Kuga--Satake varieties and the Hodge conjecture},
in: The Arithmetic and Geometry of Algebraic Cycles (Banff, 1998),
NATO Sci. Ser. C Math. Phys. Sci. \textbf{548},
Kluwer, Dordrecht, 2000, pp.~51--82.
\href{https://doi.org/10.1007/978-94-011-4098-0_3}{\texttt{doi:10.1007/978-94-011-4098-0\_3}}.

\bibitem[vG08]{vG08}
B. van Geemen,
\emph{Real multiplication on K3 surfaces and Kuga Satake varieties},
Michigan Math. J. \textbf{56} (2008), no.~2, 375--399.
\href{https://doi.org/10.1307/mmj/1224783519}{\texttt{doi:10.1307/mmj/1224783519}}.

\bibitem[vGR26]{vGR26}
B. van Geemen and A. Rapagnetta,
\emph{Hyperk\"ahler sixfolds, abelian fourfolds of Weil type and a Hodge
class},
preprint, 2026.
\href{https://arxiv.org/abs/2607.18341}{\texttt{arXiv:2607.18341}}.

\bibitem[vGS25]{vGS25}
B. van Geemen and M. Sch\"utt,
\emph{On families of K3 surfaces with real multiplication},
Forum Math. Sigma \textbf{13} (2025), Paper No. e2.
\href{https://doi.org/10.1017/fms.2024.146}{\texttt{doi:10.1017/fms.2024.146}}.

\bibitem[Voi02]{Voi02}
C. Voisin,
\emph{Hodge Theory and Complex Algebraic Geometry I},
Cambridge Studies in Advanced Mathematics \textbf{76},
Cambridge University Press, Cambridge, 2002.
\href{https://doi.org/10.1017/CBO9780511615344}{\texttt{doi:10.1017/CBO9780511615344}}.

\bibitem[Voi02b]{Voi02b}
C. Voisin,
\emph{A counterexample to the Hodge conjecture extended to K\"ahler
varieties},
Int. Math. Res. Not. \textbf{2002}, no.~20, 1057--1075.
\href{https://doi.org/10.1155/S1073792802111135}{\texttt{doi:10.1155/S1073792802111135}}.
\href{https://arxiv.org/abs/math/0112247}{\texttt{arXiv:math/0112247}}.

\bibitem[Voi03]{Voi03}
C. Voisin,
\emph{Hodge Theory and Complex Algebraic Geometry II},
Cambridge Studies in Advanced Mathematics \textbf{77},
Cambridge University Press, Cambridge, 2003.
\href{https://doi.org/10.1017/CBO9780511615177}{\texttt{doi:10.1017/CBO9780511615177}}.

\bibitem[Voi26]{Voi26}
C. Voisin,
\emph{The status of the Hodge conjecture},
guest post on the weblog of T. Tao, 12 September 2026.
\href{https://terrytao.wordpress.com/2026/09/12/the-status-of-the-hodge-conjecture/}%
{\texttt{terrytao.wordpress.com/2026/09/12/the-status-of-the-hodge-conjecture}}.

\bibitem[Wei77]{Weil77}
A. Weil,
\emph{Abelian varieties and the Hodge ring},
in: \OE uvres Scientifiques, Collected Papers, Vol.~III,
Springer, New York, 1979, pp.~421--429.

\bibitem[Xu15]{Xu15}
Z. Xu,
\emph{Algebraic cycles on a generalized Kummer variety},
Int. Math. Res. Not. IMRN \textbf{2018}, no.~3, 932--948.
\href{https://doi.org/10.1093/imrn/rnw266}{\texttt{doi:10.1093/imrn/rnw266}}.
\href{https://arxiv.org/abs/1506.04297}{\texttt{arXiv:1506.04297}}.

\bibitem[Zar83]{Zar83}
Yu. G. Zarhin,
\emph{Hodge groups of K3 surfaces},
J. Reine Angew. Math. \textbf{341} (1983), 193--220.
\href{https://doi.org/10.1515/crll.1983.341.193}{\texttt{doi:10.1515/crll.1983.341.193}}.

\bibitem[Zhu18]{Zhu18}
Y. Zhu,
\emph{K3 surfaces associated to abelian fourfolds of Mumford's type},
preprint, arXiv:1812.06882 (2018).
\href{https://doi.org/10.48550/arXiv.1812.06882}{\texttt{doi:10.48550/arXiv.1812.06882}}.

\bibitem[Zuc77]{Zuc77}
S. Zucker,
\emph{The Hodge conjecture for cubic fourfolds},
Compositio Math. \textbf{34} (1977), no.~2, 199--209.

\bibitem[Zuc79]{Zuc79}
S. Zucker,
\emph{Hodge theory with degenerating coefficients: $L_{2}$ cohomology in the
Poincar\'e metric},
Ann. of Math. (2) \textbf{109} (1979), no.~3, 415--476.
\href{https://doi.org/10.2307/1971221}{\texttt{doi:10.2307/1971221}}.

\end{thebibliography}
